\begin{table}[ht]
\centering
\caption{Sensitivity of the voting-baseline-rescaled high-value DDD to
the clip range applied to the per-tag pre/post mean-absolute-score
ratio. The four clip ranges return \emph{numerically identical}
estimates because every tag's ratio already lies within the tightest
band $[0.7, 2.0]$, so the clip never binds: the rescaling is applied
(these coefficients differ from the unrescaled $-0.131$, $-0.145$,
$-0.091$) but its result does not depend on the clip choice at all. In
particular the elite-tail estimate (score $\geq 5$) is significant even
with no clipping, which addresses the concern that the clip is a
researcher degree of freedom that rescues the result.}
\label{tab:clip_sensitivity}
\small
\begin{tabular}{lrrr}
\toprule
Clip range & score $\geq 1$ & score $\geq 2$ & score $\geq 5$ \\
\midrule
No clip            & $-0.144^{*}$ & $-0.179^{*}$ & $-0.098^{*}$ \\
$[0.3, 5.0]$       & $-0.144^{*}$ & $-0.179^{*}$ & $-0.098^{*}$ \\
$[0.5, 3.0]$ (main) & $-0.144^{*}$ & $-0.179^{*}$ & $-0.098^{*}$ \\
$[0.7, 2.0]$       & $-0.144^{*}$ & $-0.179^{*}$ & $-0.098^{*}$ \\
\bottomrule
\end{tabular}
\\[2pt]\footnotesize $^{*}$ $p<0.05$ (tag-clustered). All tag ratios lie
within $[0.7, 2.0]$, hence the rows coincide.
\end{table}
